\documentclass[10pt,letterpaper]{article}
\usepackage{cornell-notes}

\title{Clause 19.02.10: Class Overflow Error}
\author{}
\date{}
\setDocTitle{Clause 19.02.10: Class Overflow Error}
\setDocSubtitle{C++ 2024}
\setDocOwner{C++ 2024 Cornell Notes}
\setDocDate{\today}
\setCornellCollection{C++ 2024 Cornell Notes}
\setCornellUnitType{Clause}
\setCornellUnitNumber{19.02.10}
\setCornellUnitTitle{Class Overflow Error}
\hypersetup{pdftitle={Clause 19.02.10: Class Overflow Error},pdfsubject={Cornell notes},pdfauthor={}}

\begin{document}
\maketitle
\begin{CornellOverviewBox}{Overview}
\textbf{Classification:} Standard library - Normative\\[2pt]
\textbf{Learning objective:} Understand the rules, terminology, and semantic consequences associated with Class overflow\_error.
\end{CornellOverviewBox}\begin{CornellSummaryBox}{Summary}
{\sffamily\bfseries\color{Accent} Summary}\par\vspace{3pt}Section 19.2.10 focuses on Class overflow\_error. Apply the specified interface together with its mandates, effects, result conditions, exception behavior, and complexity constraints. When applying it, preserve the distinction between mandatory requirements, recommendations, permissions, and behavior deliberately left undefined, unspecified, or implementation-defined.
\end{CornellSummaryBox}

\end{document}
